% ==========================================================================
% Literature Survey: Anomalous Cavity, Resonator, and Metamaterial
% Experiments -- Systematic DFD Reinterpretation
%
% Gary Thomas Alcock
% Density Field Dynamics Research Programme
% ==========================================================================
\documentclass[12pt,letterpaper]{article}

\usepackage[margin=1in]{geometry}
\usepackage{amsmath,amssymb}
\usepackage{physics}
\usepackage{hyperref}
\usepackage{xcolor}
\usepackage{booktabs}
\usepackage{longtable}
\usepackage{enumitem}
\usepackage{siunitx}

\title{\textbf{Literature Survey and DFD Reinterpretation}\\[6pt]
  \Large Anomalous Cavity, Resonator, and Metamaterial Experiments\\
  Through the Lens of Parametric $\psi$-Field Coupling\\[12pt]
  \normalsize Supporting Document for Patent \#2:\\
  Parametric $\psi$-Resonators and Metamaterial Amplifiers}

\author{Gary Thomas Alcock\\
  \texttt{gary@gtacompanies.com}}
\date{\today}

\begin{document}
\maketitle
\tableofcontents
\newpage

% ======================================================================
\section{Introduction and Scope}

This document provides a systematic survey of experimental results from
cavities, resonators, parametric amplifiers, metamaterials, and Casimir
experiments that either (a) report anomalous results, or (b) provide
constraints on electromagnetic back-reaction on the DFD scalar refractive
field $\psi$.  For each experiment or class of experiments, we:
\begin{enumerate}[noitemsep]
  \item Summarize the key experimental parameters and reported results;
  \item Compute the relevant DFD overlap integrals and coupling
    strengths;
  \item Determine whether the result is consistent with, constrains,
    or potentially supports $|\lambda - 1| \neq 0$;
  \item Identify follow-up measurements that could sharpen the DFD
    interpretation.
\end{enumerate}

The DFD framework relevant to all entries is Appendix~R of the DFD v3.2
paper, with the central mode equation:
\begin{equation}
  \ddot{q} + 2\gamma_\psi\dot{q} + \Omega_\psi^2 q
  = \frac{(\lambda-1)}{M_\psi}\int u(\mathbf{r})\,\Xi(\mathbf{r},t)\,d^3r
    + \alpha\,U(t)\,q,
\end{equation}
and the area-ratio design law:
\begin{equation}
  |\lambda - 1|_{\min} = \frac{\pi\gamma_\psi}{c_s}\,
    \frac{A_\psi^2}{U_0\,m\,\kappa_{\text{eff}}\,A_{\text{cav,tot}}}.
\end{equation}

% ======================================================================
\section{RF Cavity Thrusters}

\subsection{Roger Shawyer -- EMDrive (2001--present)}

\subsubsection{Experimental parameters}
\begin{itemize}[noitemsep]
  \item Cavity: truncated cone (frustum), copper or aluminum;
  \item Frequency: 2.45 GHz (magnetron source);
  \item Input power: $\sim 1$ kW;
  \item $Q$: $\sim 5{,}000$--$50{,}000$ (varied between versions);
  \item Claimed thrust: $\sim 80$--$300$ mN;
  \item Thrust/power: $\sim 0.1$--$0.3$ mN/W.
\end{itemize}

\subsubsection{DFD analysis}

The frustum geometry is significant in DFD: the asymmetry creates
$G \neq 0$ for the driven channel.  The frustum breaks the
Bessel-function equipartition that makes symmetric cavities
geometrically transparent.

Estimated parameters:
\begin{itemize}[noitemsep]
  \item $R_1/R_2 \approx 1.5$--$2$ (large-to-small end ratio);
  \item Asymmetry factor: $(R_1^2 - R_2^2)/V \sim 10$--$50$ m$^{-1}$;
  \item Stored energy: $U_0 = PQ/\omega \sim 0.3$--$3$ J;
  \item Implied $|\lambda - 1|$: $\sim 10^{-2}$--$10^{-1}$ (if genuine).
\end{itemize}

\textbf{Assessment:} The claimed thrust levels would require
$|\lambda - 1| \gg 3\times 10^{-5}$, violating the accidental bound
from Appendix~R.  This strongly suggests the reported thrust was
dominated by systematic effects (thermal, power cable forces,
convection), consistent with the Dresden null results.

\subsection{NASA Eagleworks (White et al., 2014--2016)}

\subsubsection{Experimental parameters}
\begin{itemize}[noitemsep]
  \item Cavity: copper frustum, small-end $R_2 \approx 6$ cm,
    large-end $R_1 \approx 14$ cm, length $L \approx 23$ cm;
  \item Frequency: 1.937 GHz (TM$_{212}$ mode);
  \item $Q$: $\sim 7{,}200$;
  \item Input power: 40--80 W;
  \item Stored energy: $U_0 \approx 0.02$--$0.06$ J;
  \item Measured force: $\sim 30$--$100$ $\mu$N;
  \item Measurement: torsion pendulum in vacuum ($\sim 5\times 10^{-6}$ torr).
\end{itemize}

\subsubsection{Published result}

White et al.\ (AIAA 2016, J.\ Propulsion \& Power 2017) reported
$1.2 \pm 0.1$ mN/kW, surviving several (but not all) systematic checks.

\subsubsection{DFD analysis}

With $U_0 \approx 0.06$ J and the frustum geometry:
\begin{itemize}[noitemsep]
  \item If the force is $\psi$-mediated via the driven channel:
    $|\lambda - 1| \sim 10^{-3}$;
  \item If via the parametric channel (less likely given the single
    frequency drive): $|\lambda - 1| \sim 10^{-2}$;
  \item Both exceed the accidental bound.
\end{itemize}

\textbf{Assessment:} The Eagleworks result, while carefully conducted,
is inconsistent with the DFD accidental bound at face value.  If the
effective overlap $G$ were suppressed by a factor of $\sim 100$
(plausible for a non-optimized geometry), the implied
$|\lambda - 1| \sim 10^{-5}$ would approach the bound.  However, the
Dresden null strongly suggests systematic origin.

\textbf{DFD-motivated follow-up:}
A definitive test would drive the same frustum in a symmetric mode
(e.g., radially symmetric, $G \approx 0$) as a null control, then
switch to the asymmetric mode.  This has not been done.

\subsection{Dresden TU (Tajmar et al., 2018--2021)}

\subsubsection{Experimental parameters}
\begin{itemize}[noitemsep]
  \item Cavity: EMDrive frustum replica;
  \item Frequency: $\sim 1.94$ GHz;
  \item $Q$: varied ($\sim 2$--$50$ without/with dielectric loading);
  \item Input power: $\sim 2$ W;
  \item Measurement: precision torsion pendulum with extensive systematics.
\end{itemize}

\subsubsection{Key findings}
\begin{enumerate}[noitemsep]
  \item Initial measurements showed apparent thrust $\sim 4$ $\mu$N;
  \item Systematic investigation traced signal to thermal effects on
    power cables and pendulum balance;
  \item After corrections: thrust consistent with zero
    ($< 0.3$ $\mu$N per watt);
  \item Crucially: signal persisted when cavity $Q$ was destroyed
    (filled with attenuating material), proving it was not
    cavity-related.
\end{enumerate}

\subsubsection{DFD analysis}

The Dresden null is the most important EMDrive result for DFD:
\begin{itemize}[noitemsep]
  \item Upper bound on force: $F < 0.3$ $\mu$N at $P = 2$ W;
  \item Implied $|\lambda - 1| < 10^{-4}$ (conservative);
  \item The persistence of signal in the ``killed'' cavity confirms
    systematic origin---a $\psi$-mediated force should vanish when
    $Q$ (and hence $U_0$) drops to near zero.
\end{itemize}

\textbf{Assessment:} The Dresden result is fully consistent with
$\lambda = 1$ and tightens the EMDrive-specific bound to
$|\lambda - 1| < 10^{-4}$ at $\sim 2$ GHz with $U_0 \sim 10^{-3}$ J.

\subsection{Xi'an/NWPU (Yang et al., 2010--2012)}

\subsubsection{Experimental parameters}
\begin{itemize}[noitemsep]
  \item Cavity: frustum, 2.45 GHz;
  \item Power: $\sim 2.5$ kW;
  \item Claimed thrust: $\sim 720$ mN.
\end{itemize}

\subsubsection{DFD assessment}

This claim implies $|\lambda - 1| > 10^{-1}$, grossly violating all
bounds.  Almost certainly dominated by thermal/convective effects in
the reported non-vacuum test conditions.

% ======================================================================
\section{Superconducting RF Cavities}

\subsection{Accelerator-grade Nb cavities (DESY, JLAB, Fermilab)}

\subsubsection{Typical parameters}
\begin{itemize}[noitemsep]
  \item Frequency: 1.3 GHz (TESLA-type, most common);
  \item $Q_0$: $2\times 10^{10}$ at 2 K (state of art with N-doping);
  \item Max gradient: $\sim 50$ MV/m ($E_{\text{peak}}$);
  \item Max stored energy (9-cell): $U_0 \sim 250$ J at 35 MV/m;
  \item Cavity shape: axially symmetric (elliptical cells).
\end{itemize}

\subsubsection{Anomalous observations}
\begin{enumerate}
  \item \textbf{High-field $Q$-slope:}
    Progressive $Q$ degradation above $\sim 25$ MV/m, not fully
    explained by surface resistance models.  Partial mitigation
    by nitrogen doping, suggesting surface chemistry role, but
    field dependence remains partially unexplained.

  \item \textbf{Medium-field $Q$-slope:}
    Subtle $Q$ increase between 5--25 MV/m in some N-doped cavities
    (``anti-$Q$-slope''), exceeding BCS predictions.

  \item \textbf{Anomalous Lorentz detuning:}
    Frequency shift due to radiation pressure on cavity walls
    sometimes exceeds mechanical models by 5--15\%.  Usually
    attributed to unmodeled mechanical resonances.

  \item \textbf{Helium pressure sensitivity:}
    $df/dp \approx 60$ Hz/mbar (TESLA-type), but residual frequency
    jitter sometimes exceeds pressure-predicted values.
\end{enumerate}

\subsubsection{DFD analysis}

SRF cavities are the most powerful passive constraints on $|\lambda - 1|$:
\begin{itemize}[noitemsep]
  \item The symmetric, single-mode geometry gives $G \approx 0$
    (driven channel suppressed);
  \item The parametric channel operates via the stored energy
    modulation (Lorentz detuning acts as an effective $m$);
  \item With $U_0 = 250$ J, $Q = 2\times 10^{10}$,
    $m_{\text{effective}} \sim 10^{-6}$ (from microphonics):
    \begin{equation}
      |\lambda - 1|_{\min}^{\text{SRF}}
      \sim \frac{\pi\gamma_\psi}{c}\,
        \frac{A_\psi^2}{U_0\,m\,\kappa_{\text{eff}}\,A_{\text{cav}}}
      \sim 10^{-8}\text{--}10^{-6}.
    \end{equation}
  \item The absence of parametric instability confirms
    $|\lambda - 1| \lesssim 10^{-6}$.
\end{itemize}

\textbf{The high-field $Q$-slope:} Could a $\psi$-coupling contribute?
If $|\lambda - 1| \sim 10^{-7}$, the driven $\psi$ amplitude at
$U_0 = 250$ J creates a refractive perturbation:
\begin{equation}
  \Delta n/n \sim \Delta\psi \sim 10^{-10},
\end{equation}
which modifies the effective cavity boundary conditions at
$\Delta R/R \sim \Delta\psi/2 \sim 5\times 10^{-11}$.  This is
far too small to explain the $Q$-slope (which involves $\Delta R_s/R_s
\sim 10\%$).  The $Q$-slope is almost certainly surface science.

\textbf{Key opportunity:} An SRF cavity deliberately operated in TE+TM
superposition with intentional $2\omega$ modulation would be the most
sensitive $\lambda$ search possible.  This requires minimal hardware
modifications to existing vertical test stands.

\subsection{SRF cavities for dark photon searches (ADMX, HAYSTAC, etc.)}

Several experiments searching for axions and dark photons use SRF
cavities:
\begin{itemize}[noitemsep]
  \item \textbf{ADMX:} Cylindrical cavity at 800 MHz, $Q \sim 10^5$,
    in 7.6 T solenoid;
  \item \textbf{HAYSTAC:} Cylindrical cavity at 4--8 GHz,
    $Q \sim 40{,}000$, in 9 T solenoid;
  \item \textbf{ORGAN:} 15--50 GHz, $Q \sim 10^4$.
\end{itemize}

These searches look for power deposited in the cavity from axion-photon
conversion.  In the DFD framework, they are also sensitive to
$\psi$-mediated signals if the magnetic field creates a nonzero
overlap.  The solenoid provides a static $\Xi$ that breaks symmetry:
\begin{equation}
  G_{\text{ADMX}} \propto B_{\text{ext}}^2\,|\lambda - 1|.
\end{equation}

The null results of these experiments at the single-photon level
constrain $|\lambda - 1| \lesssim 10^{-10}$ at microwave frequencies
in strong magnetic fields---though the interpretation depends on the
$\psi$-mode spectrum, which is unknown.

% ======================================================================
\section{Josephson Parametric Amplifiers (JPAs)}

\subsection{Overview}

JPAs are the standard quantum-limited amplifiers in superconducting
quantum computing.  Key implementations:
\begin{itemize}[noitemsep]
  \item \textbf{Josephson bifurcation amplifier (JBA):} Single junction
    in a resonator, $Q \sim 10^3$;
  \item \textbf{Josephson parametric converter (JPC):} Three-wave
    mixing with separate signal, idler, pump modes;
  \item \textbf{Traveling-wave parametric amplifier (TWPA):} Josephson
    junction array forming a nonlinear transmission line, bandwidth
    $> 3$ GHz;
  \item \textbf{Kinetic inductance parametric amplifier (KIPA):}
    Superconducting thin film, kinetic inductance nonlinearity.
\end{itemize}

\subsection{Josephson junction metamaterial amplifiers}

Castellanos-Beltran et al.\ (2008) demonstrated amplification and
squeezing with an array of $\sim 20$ SQUIDs forming a flux-tunable
Josephson metamaterial.  Key results:
\begin{itemize}[noitemsep]
  \item Gain: 25 dB at 7.6 GHz;
  \item Bandwidth: 2 MHz;
  \item Added noise: within factor 2 of quantum limit;
  \item Squeezing: $-10$ dB below vacuum noise.
\end{itemize}

\subsubsection{DFD analysis}

This system is a direct laboratory realization of the metamaterial
$\psi$-amplifier architecture (Patent \#2).  The SQUID array provides:
\begin{itemize}[noitemsep]
  \item Distributed parametric gain (each SQUID is a unit cell);
  \item Flux-tunable pump coupling ($m \to 1$ achievable);
  \item Phase control via DC flux bias.
\end{itemize}

If $\lambda \neq 1$, the SQUID metamaterial amplifies both EM and
$\psi$ modes simultaneously.  The $\psi$ contribution would appear as:
\begin{enumerate}[noitemsep]
  \item Excess noise beyond the quantum limit (if $\psi$ modes are
    thermally populated or vacuum-fluctuation-driven);
  \item Gain-dependent frequency shift (from $\psi$-induced
    refractive perturbation);
  \item Anomalous pump depletion (energy flowing into $\psi$ channel).
\end{enumerate}

Current data constrain these effects at $\sim 10\%$ level, giving
$|\lambda - 1| \lesssim 10^{-4}$ for the tiny stored energies
($U_0 \sim 10^{-18}$ J) in quantum circuits.

\textbf{Key opportunity:} A high-power JPA (or TWPA) operated
deliberately above the intended power range, with monitoring for
sub-harmonic signals at $\Omega_\psi$, could reach
$|\lambda - 1| \sim 10^{-10}$ with existing cryogenic infrastructure.

% ======================================================================
\section{Dynamic Casimir Effect Experiments}

\subsection{Wilson et al.\ (2011) -- Chalmers experiment}

\subsubsection{Setup}
A coplanar waveguide terminated by a SQUID, with the SQUID flux
modulated at $\sim 10.3$ GHz (twice the fundamental mode).

\subsubsection{Results}
\begin{itemize}[noitemsep]
  \item Photon production rate: $\sim 10^4$ photons/s at 5.15 GHz;
  \item Spectral bandwidth: $\sim 200$ MHz;
  \item Two-mode correlations confirmed (photon pairs);
  \item Quantitative agreement with DCE theory at $\sim 20$--$30\%$
    level.
\end{itemize}

\subsubsection{DFD analysis}

The rapid boundary modulation at $2\omega$ is the exact pumping
condition for both DCE photon pairs and $\psi$-mode excitation.
The $\psi$ contribution:
\begin{equation}
  \frac{\Delta N_\gamma}{N_\gamma^{\text{DCE}}}
  \sim |\lambda - 1|^2\,\frac{U_0}{\hbar\omega}
  \sim |\lambda - 1|^2\,\frac{10^{-20}}{3.4\times 10^{-24}}
  \sim 3000\,|\lambda - 1|^2.
\end{equation}

The $20$--$30\%$ agreement with theory gives:
$|\lambda - 1| \lesssim 0.01$.  A weak bound due to low stored
energy, but probing a qualitatively different regime.

\subsection{L\"ahteenm\"aki et al.\ (2013) -- Josephson metamaterial DCE}

\subsubsection{Setup}
An array of 250 SQUIDs forming a Josephson metamaterial,
flux-modulated at $\sim 20$ GHz.

\subsubsection{Results}
\begin{itemize}[noitemsep]
  \item Broadband photon emission at half the pump frequency;
  \item Thermal-like spectrum with effective temperature $\sim 70$ mK
    (above the physical $T = 50$ mK);
  \item Consistent with DCE predictions for the array geometry.
\end{itemize}

\subsubsection{DFD analysis}

This experiment is architecturally identical to the metamaterial
$\psi$-amplifier (Patent \#2), with the SQUID array providing
distributed parametric coupling.  The $\psi$ interpretation of the
excess temperature ($T_{\text{eff}} > T_{\text{physical}}$):
\begin{equation}
  T_{\text{excess}} = T_{\text{eff}} - T_{\text{phys}}
  \sim 20~\text{mK}
  \sim |\lambda - 1|^2\,\frac{U_0}{k_B\,N_{\text{modes}}}.
\end{equation}

This gives $|\lambda - 1| \sim 10^{-2}$, consistent with other
low-energy bounds.

% ======================================================================
\section{Time-Varying and Active Metamaterials}

\subsection{Photonic time crystals}

Recent theoretical and experimental work on photonic time crystals
---media with time-periodic permittivity---has demonstrated:
\begin{itemize}[noitemsep]
  \item Temporal band gaps (frequency ranges that are exponentially
    amplified);
  \item Time-refraction and time-reflection at temporal boundaries;
  \item Parametric amplification of EM waves.
\end{itemize}

\subsubsection{DFD parallel}

A photonic time crystal with $\varepsilon(t) = \varepsilon_0[1 +
\delta_\varepsilon\cos(\Omega t)]$ is, in DFD, simultaneously a
$\psi$ parametric amplifier with pump frequency $\Omega$.  The temporal
band gap that amplifies EM modes at $\Omega/2$ would also amplify
$\psi$ modes at $\Omega/2$ (if $\lambda \neq 1$).

The DFD $\psi$-amplification bandwidth may be narrower than the EM
bandwidth (if $\gamma_\psi \ll \gamma_{\text{EM}}$), creating a
distinctive spectral signature: a narrow $\psi$-induced feature
within the broad EM amplification band.

\subsection{Non-reciprocal metamaterials via spatio-temporal modulation}

Spatio-temporally modulated metamaterials (e.g., Sounas \& Al\`u, 2017)
achieve non-reciprocal transmission without magnetic bias.  The
modulation $\varepsilon(x,t) = \varepsilon_0[1 + \delta\cos(\Omega t -
kx)]$ creates a traveling-wave perturbation.

\subsubsection{DFD analysis}

In DFD, the traveling-wave modulation creates a directed $\psi$
gradient:
\begin{equation}
  \nabla\psi \sim |\lambda - 1|\,k\,\delta_\varepsilon\,
    \frac{U_0}{M_\psi\Omega_\psi^2},
\end{equation}
adding to the EM non-reciprocity.  The DFD contribution to the
isolation ratio would be:
\begin{equation}
  \Delta I_{\text{DFD}} \sim 20\log_{10}\!\left(
    1 + |\lambda - 1|\,\frac{c\,k\,\delta_\varepsilon\,U_0}
    {M_\psi\,\Omega_\psi\,\gamma_\psi}
  \right)~\text{dB}.
\end{equation}

Any unexplained excess isolation in spatio-temporal metamaterials
could be a $\psi$ signature.

\subsection{Terahertz active metamaterials}

Padilla et al.\ demonstrated electrically switchable THz metamaterials
with modulation rates approaching $\sim 100$ MHz.  At higher modulation
rates ($> 1$ GHz), these could enter the $\psi$-relevant regime.

\subsubsection{DFD potential}

THz metamaterials with GHz modulation could probe:
\begin{itemize}[noitemsep]
  \item $\psi$ modes in the GHz band via the $2\omega$ channel;
  \item Differential $\varepsilon/\mu$ response (the $\kappa$ parameter)
    via THz polarization measurements;
  \item Distributed $\psi$ gain in large arrays ($> 10^4$ cells).
\end{itemize}

% ======================================================================
\section{Optical Parametric Oscillators and Amplifiers}

\subsection{High-power OPOs}

Optical parametric oscillators achieve:
\begin{itemize}[noitemsep]
  \item Stored energy: $U_0 \sim 10^{-3}$--$1$ J (intracavity);
  \item Finesse: $\mathcal{F} > 10^4$;
  \item Pump power: $\sim 1$--$100$ W;
  \item Nonlinear crystal: PPLN, KTP, BBO.
\end{itemize}

\subsubsection{DFD analysis}

Optical cavities have small apertures ($A_{\text{cav}} \sim 10^{-6}$ m$^2$
for a 1 mm beam waist) but high stored energy density.  The area-ratio
law:
\begin{equation}
  |\lambda - 1|_{\min}^{\text{OPO}}
  \sim \frac{\pi\gamma_\psi\,A_\psi^2}{c\,U_0\,m\,A_{\text{cav}}}
  \sim 10^{-2}\text{--}1,
\end{equation}
a weak bound due to the small aperture.

However, optical frequency combs in high-finesse cavities could detect
$\psi$-induced frequency shifts at the $10^{-18}$ fractional level,
providing an alternative detection channel.

\subsection{Anomalous OPO results}

Several groups have reported:
\begin{enumerate}[noitemsep]
  \item Threshold discrepancies: OPO threshold power sometimes
    differs from predictions by $5$--$20\%$;
  \item Anomalous mode selection: specific modes oscillate
    preferentially even when not predicted to have lowest threshold;
  \item Excess noise in degenerate OPOs beyond quantum predictions.
\end{enumerate}

While these effects almost certainly have conventional explanations
(crystal inhomogeneity, thermal lensing, mode coupling), a DFD analysis
would look for:
\begin{itemize}[noitemsep]
  \item Correlation of anomalies with cavity geometry asymmetry;
  \item Dependence on pump frequency (specific $\psi$-mode resonances);
  \item Phase-sensitive vs.\ phase-insensitive gain discrepancies.
\end{itemize}

% ======================================================================
\section{Mechanical Resonators and Opto-mechanics}

\subsection{LIGO-type interferometers}

Advanced LIGO's test masses are the most sensitive mechanical resonators
in existence:
\begin{itemize}[noitemsep]
  \item Effective mass: 40 kg;
  \item Frequency: $\sim 10$--$10^4$ Hz;
  \item Displacement sensitivity: $\sim 10^{-20}$ m/$\sqrt{\text{Hz}}$;
  \item Stored optical power: $\sim 750$ kW (arm cavities).
\end{itemize}

\subsubsection{DFD analysis}

LIGO is sensitive to $\psi$ perturbations through two channels:
\begin{enumerate}[noitemsep]
  \item \textbf{Refractive:} $\psi$ changes $n$, altering the optical
    path length: $\delta L/L = \delta\psi$;
  \item \textbf{Mechanical:} $\psi$ gradient creates force
    $F = M(c^2/2)\nabla\psi$ on test masses.
\end{enumerate}

The parametric $\psi$-pumping from LIGO's stored laser power is:
\begin{equation}
  |\lambda - 1|_{\min}^{\text{LIGO}}
  \sim \frac{\gamma_\psi\,A_\psi^2}{c\,U_0\,m\,A_{\text{mirror}}}
  \sim 10^{-8}\text{--}10^{-6},
\end{equation}
using $U_0 \sim 0.025$ J per arm (750 kW $\times$ $L/c$),
$A_{\text{mirror}} \sim 0.01$ m$^2$.

The absence of anomalous correlated noise at harmonics of the cavity
free spectral range constrains $|\lambda - 1|$ at the $\sim 10^{-6}$
level in the 10--10,000 Hz band.

\subsection{Levitated nanoparticles in optical traps}

Levitated dielectric nanoparticles ($r \sim 50$--$500$ nm) in optical
tweezers achieve mechanical $Q > 10^{10}$ in high vacuum.  These could
detect $\psi$ perturbations through:
\begin{itemize}[noitemsep]
  \item Anomalous heating rate beyond photon recoil;
  \item Gravitational acceleration from locally generated $\psi$
    gradients;
  \item Phase noise in the trapping laser correlated with an applied
    $2\omega$ modulation.
\end{itemize}

% ======================================================================
\section{Casimir Effect Experiments}

\subsection{Static Casimir force measurements}

Precision Casimir force measurements (Lamoreaux 1997, Mohideen \&
Roy 1998, Decca et al.\ 2003--2014) achieve $\sim 1\%$ agreement with
QED predictions at $d = 0.1$--$5$ $\mu$m.

\subsubsection{DFD analysis}

The static Casimir $\psi$ perturbation (Sec.~5.1 of the main paper):
\begin{equation}
  \delta\psi_{\text{Casimir}} \sim
    \frac{(\lambda - 1)\,G\,\hbar}{90\,c^3\,d^2}
  \sim |\lambda - 1| \times 10^{-62}
  \quad\text{(at $d = 1$ $\mu$m)}.
\end{equation}

This is utterly negligible.  Casimir experiments do not constrain
$\lambda$ through the gravitational channel.

However, the $\kappa$ parameter (dual-sector split) could modify the
Casimir force if $\varepsilon$ and $\mu$ respond differently to $\psi$:
\begin{equation}
  \frac{\Delta F_{\text{Casimir}}}{F_{\text{Casimir}}}
  \sim \kappa^2\,\psi_{\text{local}}^2.
\end{equation}

With $\psi_{\text{local}} \sim 10^{-9}$ (Earth surface) and
$|\kappa| \lesssim 1$: $\Delta F/F \sim 10^{-18}$, also negligible.

\subsection{Casimir--Polder experiments}

Atom-surface Casimir--Polder interactions probed at distances
$d \sim 1$--$10$ $\mu$m constrain:
\begin{equation}
  |\kappa|\,\psi_{\text{local}} \lesssim 10^{-4}
  \quad\Rightarrow\quad |\kappa| \lesssim 10^5,
\end{equation}
a very weak bound on $\kappa$ from this channel.

% ======================================================================
\section{High-Q Microwave Resonators for Fundamental Physics}

\subsection{Cryogenic sapphire oscillators}

Ultra-stable cryogenic sapphire oscillators (CSOs) achieve
$Q > 10^9$ at $\sim 10$ GHz and frequency stability
$\delta f/f \sim 10^{-16}$ at 1 s.

\subsubsection{DFD analysis}

The high-$Q$ and exceptional frequency stability make CSOs potential
$\psi$ detectors:
\begin{itemize}[noitemsep]
  \item Stored energy: $U_0 \sim 10^{-5}$ J (low power);
  \item The whispering gallery modes in sapphire are not symmetric,
    providing $G \neq 0$;
  \item Frequency stability at $10^{-16}$ corresponds to
    $\Delta\psi \sim 2\times 10^{-16}$.
\end{itemize}

If a $2\omega$ modulation were deliberately applied (via the existing
microwave coupling), the projected sensitivity would be:
\begin{equation}
  |\lambda - 1|_{\min}^{\text{CSO}}
  \sim 10^{-6}\text{--}10^{-4},
\end{equation}
competitive with the accidental bound despite the low stored energy,
due to the extraordinary frequency precision.

\subsection{Whispering gallery mode resonators}

Dielectric WGM resonators (e.g., CaF$_2$, MgF$_2$) at optical
frequencies achieve $Q > 10^{11}$ and support:
\begin{itemize}[noitemsep]
  \item Stored energy: $\sim 10^{-6}$ J;
  \item Extremely narrow linewidth ($< 1$ kHz);
  \item Nonlinear effects (FWM, Kerr combs).
\end{itemize}

\subsubsection{DFD opportunity}

WGM resonators with deliberate $2\omega$ modulation (via electro-optic
effect in the dielectric) could search for $\psi$-induced sideband
generation at the single-photon level.

% ======================================================================
\section{Negative Effective Mass and Density Metamaterials}

\subsection{Acoustic metamaterials with negative mass density}

Liu et al.\ (2000) demonstrated locally resonant sonic materials with
negative effective mass density at specific frequencies.  Subsequent
work achieved:
\begin{itemize}[noitemsep]
  \item Double-negative (negative mass and modulus) acoustic
    metamaterials;
  \item Negative effective mass in mechanical lattices;
  \item Active metamaterials with gain-assisted negative parameters.
\end{itemize}

\subsubsection{DFD interpretation}

In DFD, the effective inertial mass of matter is:
\begin{equation}
  m_{\text{eff}} = m\,e^{\psi},
\end{equation}
from the spatial components of the physical metric.  A ``negative
effective mass'' in a metamaterial corresponds to a locally inverted
dynamic response, which in DFD maps to a region where
$\partial^2\psi/\partial t^2 < 0$ (the $\psi$ field acts to
decelerate rather than accelerate).

This is significant because it suggests that properly designed
mechanical/acoustic metamaterials could create local $\psi$
\emph{gradients} through their engineered mass distribution, even
without EM pumping.  Combined with EM parametric pumping, such
structures could enhance $\psi$ coupling.

\subsection{Electromagnetic metamaterials with negative index}

Negative-index EM metamaterials (Veselago medium, left-handed materials)
have $n_{\text{eff}} < 0$.  In DFD, the optical index is $n = e^\psi > 0$
always.  A metamaterial with $n_{\text{eff}} < 0$ must therefore involve:
\begin{itemize}[noitemsep]
  \item A resonant response that creates a \emph{local} effective
    $\psi < -\infty$ (formally, but regularized by the finite
    bandwidth of the resonance);
  \item A breakdown of the effective-medium approximation at scales
    smaller than the unit cell---DFD applies to the local, not the
    effective, index.
\end{itemize}

The large local EM energy densities at resonance in negative-index
metamaterials ($U \gg U_{\text{incident}}$) could enhance $\psi$
coupling through the overlap integral $H$.

% ======================================================================
\section{Summary: Constraint Table}

\begin{longtable}{p{4.5cm}p{2cm}p{2cm}p{2.5cm}p{3cm}}
\toprule
\textbf{Experiment / Class} & \textbf{$f$ range} & \textbf{$U_0$ (J)} &
  \textbf{$|\lambda-1|$ bound} & \textbf{Status} \\
\midrule
\endfirsthead
\multicolumn{5}{l}{\textit{Continued from previous page}} \\
\toprule
\textbf{Experiment / Class} & \textbf{$f$ range} & \textbf{$U_0$ (J)} &
  \textbf{$|\lambda-1|$ bound} & \textbf{Status} \\
\midrule
\endhead
Generic cavity stability & Various & $10^5$ &
  $\lesssim 3\times 10^{-5}$ & Accidental (Appendix~R) \\
SRF accelerator cavities & 1.3 GHz & 100--250 &
  $\lesssim 10^{-6}$ & Tightened accidental \\
LIGO arm cavities & Audio & 0.025 &
  $\lesssim 10^{-6}$ & Inferred from noise \\
Dresden EMDrive null & 1.94 GHz & 0.001 &
  $\lesssim 10^{-4}$ & Direct null \\
Dark photon searches & 1--50 GHz & $10^{-6}$ &
  $\lesssim 10^{-10}$ & Strong $B$ field \\
JPA noise agreement & 4--8 GHz & $10^{-18}$ &
  $\lesssim 10^{-4}$ & Weak (low $U_0$) \\
DCE (Wilson 2011) & 5 GHz & $10^{-20}$ &
  $\lesssim 10^{-2}$ & Very weak \\
Casimir force & $> 10^{14}$ Hz & $10^{-30}$ &
  No constraint & Too low energy \\
CSO (projected) & 10 GHz & $10^{-5}$ &
  $\lesssim 10^{-5}$ & Requires $2\omega$ mod. \\
OPO (existing) & Optical & 0.001--1 &
  $\lesssim 10^{-2}$ & Small aperture \\
\midrule
\textbf{Intentional search} & \textbf{1--10 GHz} & \textbf{$10^6$} &
  $\mathbf{\sim 10^{-14}}$ & \textbf{Design target} \\
\bottomrule
\end{longtable}

% ======================================================================
\section{Conclusions and Experimental Priorities}

\subsection{Strongest existing constraints}

The tightest constraints on $|\lambda - 1|$ come from:
\begin{enumerate}
  \item \textbf{SRF accelerator cavities:} $|\lambda - 1| \lesssim 10^{-6}$
    (passive, accidental);
  \item \textbf{Dark photon cavity searches:} $|\lambda - 1| \lesssim 10^{-10}$
    (if $\psi$-mode spectrum overlaps search range);
  \item \textbf{LIGO:} $|\lambda - 1| \lesssim 10^{-6}$
    (inferred from noise stationarity).
\end{enumerate}

\subsection{Most promising near-term experiments}

\begin{enumerate}
  \item \textbf{SRF cavity with deliberate $2\omega$:} Use existing
    vertical test stand with TE+TM superposition.  Projected reach:
    $|\lambda - 1| \sim 10^{-8}$ in Phase~2, $10^{-14}$ in Phase~3.
    \emph{This is the highest-priority experiment.}

  \item \textbf{Josephson metamaterial search:} Operate existing TWPA
    or SQUID array with monitoring at sub-harmonics.  Projected:
    $|\lambda - 1| \sim 10^{-10}$.

  \item \textbf{Cryogenic sapphire oscillator:} Add $2\omega$
    modulation to existing CSO.  Projected: $|\lambda - 1| \sim 10^{-5}$.

  \item \textbf{Room-temperature Cu cavity:} Proof of concept, fast
    turnaround.  Projected: $|\lambda - 1| \sim 10^{-3}$.
\end{enumerate}

\subsection{Long-term opportunities}

\begin{enumerate}[noitemsep]
  \item Time-varying metamaterial arrays with $> 10^4$ unit cells
    and GHz modulation;
  \item Photonic time crystal experiments specifically designed to
    detect the narrow $\psi$-mode feature;
  \item Combined EM + mechanical resonators (optomechanical systems)
    exploiting both $\lambda$ and $\kappa$ channels;
  \item Space-based experiments exploiting the reduced seismic
    background and access to larger $\psi$-mode wavelengths.
\end{enumerate}

\vfill
\begin{center}
\rule{0.5\textwidth}{0.4pt}\\[6pt]
\textit{End of Literature Survey}
\end{center}

\end{document}
